\chapter{Group Theory}

A \emph{group} is a set \(G\) equipped with a binary operation \(\cdot: G \times G \to G\) satisfying the following axioms:

\begin{enumerate}

    \item \emph{Associativity:} For all \(a, b, c \in G\), \((a \cdot b) \cdot c = a \cdot (b \cdot c)\).

    \item \emph{Identity element:} There exists an element \(e \in G\) such that for all \(a \in G\), \(e \cdot a = a \cdot e = a\).

    \item \emph{Inverse element:} For each \(a \in G\), there exists an element \(a^{-1} \in G\) such that \(a \cdot a^{-1} = a^{-1} \cdot a = e\).

\end{enumerate}o

\section{Symmetric Group}

Given a set \(G\). The set of all bijections from \(G\) to itself, denoted by \(S_G\), forms a group under composition of functions. 
This group is called the \emph{symmetric group} on \(G\).

For \(G = \{1, 2, \dots, n\}\), the symmetric group on \(G\) is usually denoted by \(S_n\).



